\subsection{A reduction theorem}

We preserve the notation and hypotheses of nos. 2 to 7.

LEMMA 6. Let R be a commutative ring and \(\mathfrak{p}_i\) ( \(1 \leqslant i \leqslant n\) ) distinct prime ideals of R.

\begin{enumerate}
\def\labelenumi{(\roman{enumi})}
\item
  For \(1 \leqslant i \leqslant n\) , let \(\mathbf{H}_i\) be a subset of \(\mathbf{R} / \mathfrak{p}_i\) satisfying the following condition: there exists no element \(\alpha_i \in \mathbf{R} / \mathfrak{p}_i\) such that \(\alpha_i + \mathbf{H}_i\) contains an ideal \(\neq 0\) of \(\mathbf{R} / \mathfrak{p}_i\) . Then there exists \(a \in \mathbf{R}\) such that, for \(1 \leqslant i \leqslant n\) , the canonical image of \(a\) in \(\mathbf{R} / \mathfrak{p}_i\) does not belong to \(\mathbf{H}_i\) .
\item
  If \(\operatorname{Card}(\mathbf{H}_i) < \operatorname{Card}(\mathbf{R} / \mathfrak{p}_i)\) , the \(\mathbf{H}_i\) satisfy the condition in (i).
\item
  We argue by induction on \(n\) , the case \(n = 0\) being trivial. Therefore let \(n \geqslant 1\) . We may make a permutation on the indices \(i\) and may assume that \(p_1\) is minimal among the \(p_i\) and therefore, for \(2 \leqslant i \leqslant n\) , there exists \(c_i \in p_i\) such that \(c_i \notin p_1\) . By the induction hypothesis there exists \(b \in R\) such that the canonical image of \(b\) in \(R / p_i\) does not belong to \(H_i\) for \(2 \leqslant i \leqslant n\) . For all \(x \in R\) we write \(a_{x}=b+xc_{2}c_{3}\ldots c_{n};\) as \(c_{i}\in\mathfrak{p}_{i},\) obviously \(a_{x}\equiv b\pmod{\mathfrak{p}_{i}}\) for \(2\leqslant i\leqslant n\) . It therefore suffices to prove that there exists \(x\in R\) such that the canonical image of \(a_{x}\) in \(R/\mathfrak{p}_{1}\) does not belong to \(H_{1}\) . Now, the set of canonical images of the \(a_{x}\) in \(R/\mathfrak{p}_{1}\) , where x runs through R, is just \(\beta+c\) , where \(\beta\) is the canonical image of b and c the ideal of \(R/\mathfrak{p}_{1}\) generated by the canonical image of \(c_{2}c_{3}\ldots c_{n}\) ; by virtue of the choice of the \(c_{i},c\neq0\) since \(R/\mathfrak{p}_{1}\) is an integral domain and the hypothesis on \(H_{1}\) implies the existence of an x which solves the problem.
\item
  As \(R/p_{i}\) is an integral domain, every ideal \(\neq0\) of \(R/p_{i}\) has cardinal equal to that of \(R/p_{i}\) and the same is true of any translation of an ideal by an element of \(R/p_{i}\) , whence the conclusion.
\end{enumerate}

THEOREM 6. Let M be a torsion-free finitely generated A-module. There exists a free submodule L of M such that M/L is isomorphic to an ideal of A.

We shall denote by \(n\) the rank of \(M\) (rank of \(V = M \otimes_{A} K\) over \(K\) ) and consider \(M\) as a lattice of \(V\) with respect to \(A\) . Then, for all \(p \in P(A)\) , \(M_p\) is a lattice of \(V\) with respect to \(A_p\) (no. 1, Example 6) and, as \(A_p\) is a principal ideal domain, \(M_p\) is a free \(A_p\) -module of rank \(n\) . We shall write:

\[
\mathbf {M} (\mathfrak {p}) = \mathbf {M} _ {\mathfrak {p}} / \mathfrak {p} \mathbf {M} _ {\mathfrak {p}}.
\]

We shall denote by \(k(\mathfrak{p})\) the field of fractions of \(\mathbf{A} / \mathfrak{p}\) (isomorphic to the residue field of \(\mathbf{A}_{\mathfrak{p}}\) ); \(\mathbf{M}(\mathfrak{p}) = \mathbf{M} \otimes_{\mathbf{A}} k(\mathfrak{p})\) is therefore a vector space of rank \(n\) over \(k(\mathfrak{p})\) . For all \(x \in \mathbf{M}\) , we shall denote by \(x(\mathfrak{p})\) the canonical image of \(x\) in \(\mathbf{M}(\mathfrak{p})\) .

LEMMA 7. Let \(x_{i}\) ( \(1 \leqslant i \leqslant m\) ) be linearly independent elements of \(\mathbf{M}\) (over \(\mathbf{A}\) or \(\mathbf{K}\) , which amounts to the same) and let \(\mathbf{L}\) be the sub-A-module of \(\mathbf{M}\) generated by the \(x_{i}\) . Then, for almost all \(\mathfrak{p} \in \mathbf{P}\) , the \(x_{i}(\mathfrak{p}) \in \mathbf{M}(\mathfrak{p})\) are linearly independent over \(k(\mathfrak{p})\) ; for them to be linearly independent over \(k(\mathfrak{p})\) for all \(\mathfrak{p} \in \mathbf{P}\) , it is necessary and sufficient that \(\mathbf{M}/\mathbf{L}\) be torsion-free.

Let \(x_{m+1}, \ldots, x_n\) be elements of M which, with \(x_1, \ldots, x_m\) , form a basis of V and let N be the free sub-A-module of M generated by the \(x_i (1 \leqslant i \leqslant n)\) . It follows from no. 3, Theorem 3 that \(N_p = M_p\) for almost all \(p \in P\) ; as \(x_1(p), \ldots, x_n(p)\) form a basis of \(N(p)\) over \(k(p)\) , this establishes the first assertion.

If \(\mathbf{M} / \mathbf{L}\) is torsion-free, so is \((\mathbf{M} / \mathbf{L})_{\mathfrak{p}} = \mathbf{M}_{\mathfrak{p}} / \mathbf{L}_{\mathfrak{p}}\) for all \(\mathfrak{p} \in \mathbb{P}\) (no. 1, Example 6) and, as \(\mathbf{A}_{\mathfrak{p}}\) is a principal ideal domain, \(\mathbf{M}_{\mathfrak{p}} / \mathbf{L}_{\mathfrak{p}}\) is free. Therefore \(\mathbf{M}_{\mathfrak{p}}\) is the direct sum of \(\mathbf{L}_{\mathfrak{p}}\) and a free \(\mathbf{A}_{\mathfrak{p}}\) -module E of rank \(n - m\) ; hence \(\mathbf{M}(\mathfrak{p})\) is the direct sum of \(\mathbf{L}(\mathfrak{p})\) and the vector \(k(\mathfrak{p})\) -space \(\mathrm{E} / \mathfrak{p}\mathrm{E}\) of rank \(n - m\) ; therefore \(\mathbf{L}(\mathfrak{p})\) is of rank \(m\) and, as it is generated by the \(x_{i}(\mathfrak{p})\) ( \(1 \leqslant i \leqslant m\) ), the latter are linearly independent.

Conversely, suppose that the \(x_{i}(\mathfrak{p})\) ( \(1 \leqslant i \leqslant m\) ) are linearly independent over \(k(\mathfrak{p})\) for all \(\mathfrak{p} \in \mathbb{P}\) . Then \(\mathbf{L}_{\mathfrak{p}}\) is a direct factor of \(\mathbf{M}_{\mathfrak{p}}\) for all \(\mathfrak{p}\) (Chapter II, §3, no. 2, Corollary 1 to Proposition 5) and therefore \(\mathbf{M}_{\mathfrak{p}} / \mathbf{L}_{\mathfrak{p}} = (\mathbf{M} / \mathbf{L})_{\mathfrak{p}}\) is torsion-free for all \(p \in P\) . We conclude that \(\mathrm{P} \cap \mathrm{Ass}(\mathrm{M}/\mathrm{L}) = \varnothing\) by Chapter IV, §1, no. 2, Corollary to Proposition 5. But as L is reflexive, it follows from no. 2, Proposition 7 (i) that the only prime ideal which can belong to \(\mathrm{Ass}(\mathrm{M}/\mathrm{L})\) is the ideal (0); hence M/L is torsion-free.

LEMMA 8. Suppose that the rank n of M is \(\geqslant2\) ; then there exists an element \(x \neq 0\) of M such that M/Ax is torsion-free.

Let \(y \neq 0\) be an element of M. By Lemma 7, the set Y of \(p \in P\) such that \(y(p) = 0\) is finite. If Y = \(\varnothing\) , it follows from Lemma 7, applied to the sequence \((x_i)\) consisting of the single element y, that M/Ay is torsion-free. Suppose therefore that Y ≠ \(\varnothing\) and write S = \(\bigcap_{p \in Y} (A - p)\) ; we know (no. 4, Lemma 2) that \(S^{-1}A\) is a semi-local principal ideal domain whose maximal ideals are the \(pS^{-1}A\) , where p ∈ Y, the corresponding local rings being the \(A_{p}\) . Therefore

\[
\mathrm{S} ^ {- 1} \mathrm{A} / \mathfrak {p} \mathrm{S} ^ {- 1} \mathrm{A} = k (\mathfrak {p}),
\]

whence

\[
\begin{array}{c} \mathrm{S} ^ {- 1} \mathrm{M} / \mathfrak {p} \mathrm{S} ^ {- 1} \mathrm{M} = (\mathrm{M} / \mathfrak {p} \mathrm{M}) \otimes_ {\mathrm{A}} \mathrm{S} ^ {- 1} \mathrm{A} = \mathrm{M} \otimes_ {\mathrm{A}} ((\mathrm{A} / \mathfrak {p}) \otimes_ {\mathrm{A}} \mathrm{S} ^ {- 1} \mathrm{A}) = \\ \mathrm{M} \otimes_ {\mathrm{A}} k (\mathfrak {p}) = \mathrm{M} (\mathfrak {p}) \end{array}
\]

for all \(\mathfrak{p} \in \mathbf{Y}\) . By Chapter II, § 1, no. 2, Proposition 6, there exists an element \(z / s \in \mathrm{S}^{-1}\mathbf{M}\) ( \(z \in \mathbf{M}, s \in \mathrm{S}\) ) whose canonical images in the \(\mathbf{M}(\mathfrak{p})\) for \(\mathfrak{p} \in \mathbf{Y}\) are all \(\neq 0\) . By definition of \(S\) therefore \(z(\mathfrak{p}) \neq 0\) for all \(\mathfrak{p} \in Y\) . We may further assume that \(y\) and \(z\) are linearly independent over \(K\) . For in the opposite case consider an element \(t \in M\) which is linearly independent of \(y\) (such exists since \(n \geqslant 2\) ); on the other hand take an element \(a \neq 0\) belonging to \(\bigcap_{\mathfrak{p} \in Y} \mathfrak{p}\) (which is not reduced to 0 since \(A\) is an integral domain) and write \(z' = z + at\) : clearly \(y\) and \(z'\) are linearly independent over \(K\) and \(z'(\mathfrak{p}) = z(\mathfrak{p}) \neq 0\) for all \(\mathfrak{p} \in Y\) .

Supposing therefore that \(y\) and \(z\) are linearly independent over \(\mathbf{K}\) , let \(\mathbf{Z}\) be the set of \(\mathfrak{p} \in \mathbf{P} - \mathbf{Y}\) such that \(y(\mathfrak{p})\) and \(z(\mathfrak{p})\) are linearly independent over \(k(\mathfrak{p})\) ; it follows from Lemma 7 that this set is finite. For all \(\mathfrak{p} \in \mathbf{Z}\) , we may therefore write \(z(\mathfrak{p}) = \lambda(\mathfrak{p})y(\mathfrak{p})\) where \(\lambda(\mathfrak{p}) \in k(\mathfrak{p})\) . Now, \(\operatorname{Card}(\mathrm{A}/\mathfrak{p}) \geqslant 2\) for all \(\mathfrak{p} \in \mathbf{P}\) ; it therefore follows from Lemma 6 that there exists \(b \in \mathbf{A}\) such that, for all \(\mathfrak{p} \in \mathbf{Z}\) , the canonical image of \(b\) in \(\mathrm{A}/\mathfrak{p}\) is distinct from \(\lambda(\mathfrak{p})\) . We now show that the element \(x = z - by\) solves the problem; it suffices (by virtue of Lemma 7 applied with \(m = 1\) ) to verify that \(x(\mathfrak{p}) \neq 0\) for all \(\mathfrak{p} \in \mathbf{P}\) . Now:

--- if \(\mathfrak{p} \in \mathbf{Y}\) , then \(x(\mathfrak{p}) \neq 0\) by construction;

--- if \(p \in Z\) , then \(x(p) = \mu_{p} y(p)\) where \(\mu_{p} \neq 0\) by virtue of the choice of b and hence \(x(p) \neq 0\) since \(y(p) \neq 0\) ;

- if \(\mathfrak{p} \in \mathrm{P} - (\mathrm{Y} \cup \mathrm{Z}), y(\mathfrak{p})\) and \(z(\mathfrak{p})\) are linearly independent and hence \(x(\mathfrak{p}) \neq 0\) .

Having established these lemmas, we pass to the proof of Theorem 6. We argue by induction on n, the case \(n \leqslant 1\) being trivial since M itself is then isomorphic to an ideal of A. Suppose therefore that \(n \geqslant 2\) ; by virtue of Lemma 8 there exists a free submodule \(L_{0}\) of M of rank 1 such that \(M/L_{0}\) is torsion-free; \(M/L_{0}\) is therefore of rank n - 1. By the induction hypothesis there exists a free submodule \(L_{1}\) of \(M/L_{0}\) such that \((\mathrm{M}/\mathrm{L}_{0})/\mathrm{L}_{1}\) is isomorphic to an ideal of A. Let L be the inverse image of \(L_{1}\) in M; \(L/L_{0}\) is isomorphic to \(L_{1}\) and, since \(L_{1}\) is free, L is isomorphic to \(L_{0} \oplus L_{1}\) (Algebra, Chapter II, § 1, no. 11, Proposition 21) and hence free; as M/L is isomorphic to \((\mathrm{M}/\mathrm{L}_{0})/\mathrm{L}_{1}\) , the theorem has been shown.

Remark. If M is reflexive, the same is not necessarily true of M/L (Exercise 9).

